\documentclass{standalone}
\usepackage{circuitikz}
\usetikzlibrary{calc}
\definecolor{circuitnotemark}{HTML}{FE00FE}
\begin{document}
\begin{circuitikz}[american, line width=0.8pt]
\ctikzset{bipoles/length=1.2cm}
\ctikzset{grounds/scale=1.36}
\coordinate (x1y2) at (0,-2);
\coordinate (x1y5) at (0,-8);
\coordinate (x3y2) at (4,-2);
\coordinate (x5y2) at (8,-2);
\coordinate (x7y2) at (12,-2);
\coordinate (x9y2) at (16,-2);
\coordinate (x7y1) at (12,0);
\coordinate (x9y1) at (16,0);
\coordinate (x9y5) at (16,-8);
\draw (x1y2) to[esource, l_=$\mathrm{W1}$] (x1y5); % line 3
\draw (-0.18,-5) sin ++(0.09,0.09) cos ++(0.09,-0.09) sin ++(0.09,-0.09) cos ++(0.09,0.09); % line 3
\node[ground] at (x1y5) {}; % line 4
\draw (x3y2) to[L, l_=$L_{1}$, a^=$10\,\mathrm{m}\mathrm{H}$] (x5y2); % line 5
\draw (x5y2) to[C, l_=$C_{1}$, a^=$10\,\mathrm{n}\mathrm{F}$] (x7y2); % line 6
\draw (x7y2) to[R, l_=$R_{1}$, a^=$100\,\Omega$] (x9y2); % line 7
\draw (x7y1) to[rmeter, t={$\mathrm{V}$}, l_=$\mathrm{CH2}$] (x9y1); % line 8
\draw (x1y2) -- (x3y2); % line 10
\draw (x7y1) -- (x7y2); % line 11
\draw (x9y1) -- (x9y2); % line 12
\draw (x9y2) -- (x9y5); % line 13
\draw (x9y5) -- (x1y5); % line 14
\node[circ] at (x1y5) {};
\node[circ] at (x7y2) {};
\node[circ] at (x9y2) {};
\node[anchor=south west, inner sep=2pt, yshift=4pt, circuitnotemark, font=\LARGE\bfseries] (circuittitle) at (current bounding box.north west) {X};
\path (circuittitle.south west) rectangle ++(7.955,0.853);
\node[anchor=north east, inner sep=2pt, circuitnotemark, font=\large] (circuitstamp) at ([yshift=-0.593cm]current bounding box.south east) {X};
\path (circuitstamp.north east) rectangle ++(-5.08,-0.593);
\end{circuitikz}
\end{document}
